\documentclass[11pt,a4paper]{article}
\usepackage{fontspec,amsmath,amssymb,amsthm,microtype}
\setmainfont{lmroman10-regular.otf}[BoldFont=lmroman10-bold.otf,ItalicFont=lmroman10-italic.otf,BoldItalicFont=lmroman10-bolditalic.otf]
\usepackage{xcolor}
\usepackage[most]{tcolorbox}
\usepackage{mathrsfs}
\providecommand{\lt}{<}
\providecommand{\gt}{>}
\setmonofont{DejaVuSansMono.ttf}[BoldFont=DejaVuSansMono-Bold.ttf,ItalicFont=DejaVuSansMono-Oblique.ttf,BoldItalicFont=DejaVuSansMono-BoldOblique.ttf,Scale=0.86]
\newfontfamily\leanfallback{STIXTwoMath-Regular.otf}[Scale=0.86]
\definecolor{leanrule}{gray}{0.62}
\definecolor{leanback}{gray}{0.962}
\definecolor{leantitle}{gray}{0.22}
\definecolor{leanmuted}{gray}{0.38}
\newtcolorbox{leanbox}[1]{enhanced,breakable,colback=leanback,colframe=leanrule,
  boxrule=0.4pt,arc=0.6mm,left=2.4mm,right=2.4mm,top=1.8mm,bottom=1.8mm,
  before skip=0.8em,after skip=1em,fonttitle=\small\bfseries,coltitle=leantitle,
  colbacktitle=leanback,titlerule=0.3pt,title={#1}}
\newlength{\leancharwidth}
\newenvironment{leancode}{\par\smallskip\begingroup\ttfamily\footnotesize\raggedright
  \settowidth{\leancharwidth}{M}\setlength{\parskip}{0pt}\setlength{\parindent}{0pt}
  \frenchspacing}{\par\endgroup\smallskip}
\newcommand{\leanline}[2]{\hangindent=\dimexpr#1\leancharwidth+4\leancharwidth\relax
  \hangafter=1\noindent\hspace*{#1\leancharwidth}#2\par}
\newcommand{\leansource}[1]{\par{\footnotesize\color{leanmuted}#1}\par}
\newcommand{\leanlabel}[1]{\par\smallskip\noindent{\small\bfseries #1}\par}
\usepackage[margin=25mm]{geometry}
\usepackage[unicode,hidelinks]{hyperref}
\hypersetup{pdfauthor={LangevinLean Project}}
\setlength{\emergencystretch}{3em}
\setlength{\parskip}{0.4em}
\setlength{\parindent}{0pt}
\allowdisplaybreaks[2]
\raggedbottom
\title{\texorpdfstring{A vanishing coupling distance determines the limiting law}{A vanishing coupling distance determines the limiting law}}
\author{LangevinLean Project}
\date{}
\begin{document}\maketitle
Draft mathematical exposition; these notes may contain mistakes.


Let \(H\) be a real inner product space with its Borel \(\sigma\)-algebra. The characteristic function of a finite measure \(\nu\) on \(H\) is


\[\varphi_\nu(t)=\int_He^{i\langle x,t\rangle}\,\nu(dx),\qquad t\in H.\]


If \(H\) is complete and separable (for instance \(H=\mathbb R^d\)), two finite measures with the same characteristic function are equal. This section proves a criterion for identifying the law of a limit. Suppose \(X_n\to X\) in \(L^2(P;H)\), and suppose that for each \(n\) there is a coupling \(\rho_n\) of the law of \(X_n\) with a fixed probability measure \(\pi\) such that the expected distance \(\int\|x-y\|\,\rho_n(dx,dy)\) tends to \(0\). Then the law of \(X\) is \(\pi\).




The proof uses two estimates for the characteristic test \(x\mapsto e^{i\langle x,t\rangle}\), which is bounded by \(1\) and \(\|t\|\)-Lipschitz: the difference of the characteristic functions of the two marginals of a coupling is at most \(\|t\|\) times its expected distance, and \(X\mapsto\varphi_{P_X}(t)\) is continuous on \(L^2\). So \(\varphi_{P_{X_n}}(t)\) converges both to \(\varphi_{P_X}(t)\) and to \(\varphi_\pi(t)\). The criterion is applied in the next section to Euler approximations of the Langevin equation coupled with Metropolis-adjusted chains that preserve the Gibbs measure. Only first moments of the couplings are used, and no optimal coupling is required.


\section{\texorpdfstring{A Lipschitz bound for characteristic tests}{A Lipschitz bound for characteristic tests}}


\textbf{Claim.} Let \(H\) be a real inner product space and \(t,x,y\in H\). Then


\[\bigl|e^{i\langle x,t\rangle}-e^{i\langle y,t\rangle}\bigr|\le\|t\|\,\|x-y\|.\]


\textbf{Proof.} For real numbers \(u,v\), the map \(\theta\mapsto e^{i\theta}\) has derivative \(ie^{i\theta}\) of modulus \(1\), so


\[\bigl|e^{iu}-e^{iv}\bigr|=\Bigl|\int_v^u ie^{i\theta}\,d\theta\Bigr|\le|u-v|:\qquad(1)\]


the unit circle is parametrised by a \(1\)-Lipschitz map. Apply (1) with \(u=\langle x,t\rangle\) and \(v=\langle y,t\rangle\). By linearity of the inner product in its first argument, \(u-v=\langle x-y,t\rangle\), and by the Cauchy–Schwarz inequality \(|\langle x-y,t\rangle|\le\|x-y\|\,\|t\|\). \(\square\)

\begin{leanbox}{Lean: \texttt{characteristicTest\_\allowbreak{}norm\_\allowbreak{}sub\_\allowbreak{}le}}
\leansource{\texttt{characteristicTest\_\allowbreak{}norm\_\allowbreak{}sub\_\allowbreak{}le} \quad UnadjustedLangevin/\allowbreak{}CouplingLimit.lean, lines 27–38}
\begin{leancode}
\leanline{0}{theorem\ characteristicTest\_norm\_sub\_le\ (t\ x\ y\ :\ H)\ :}
\leanline{4}{‖Complex.exp\ ((inner\ ℝ\ x\ t\ :\ ℂ)\ *\ Complex.I)\ -}
\leanline{6}{Complex.exp\ ((inner\ ℝ\ y\ t\ :\ ℂ)\ *\ Complex.I)‖\ ≤\ ‖t‖\ *\ ‖x\ -\ y‖}
\end{leancode}
\leanlabel{How the Lean proof follows the argument}
\texttt{characteristic\allowbreak{}Test\_\allowbreak{}norm\_\allowbreak{}sub\_\allowbreak{}le} is stated with \texttt{omit [Measurable\allowbreak{}Space H] [Borel\allowbreak{}Space H] [Second\allowbreak{}Countable\allowbreak{}Topology H] in}, since no measurability is involved. The characteristic test is written \texttt{Complex.\allowbreak{}exp ((inner \ensuremath{\mathbb R} x t : \ensuremath{\mathbb C}) * Complex.\allowbreak{}I)}, with the real inner product coerced to \(\mathbb C\).



\begin{enumerate}


\item \texttt{hc} is (1): Mathlib's \texttt{lipschitz\allowbreak{}With\_\allowbreak{}circle\allowbreak{}Map (0 : \ensuremath{\mathbb C}) 1} says that \(\theta\mapsto0+1\cdot e^{i\theta}\) is Lipschitz with constant \(|1|\) (\texttt{hunit} computes \(|1|=1\)); \texttt{simp only [circle\allowbreak{}Map, …]} rewrites it as the displayed inequality.

\item The \texttt{calc} block then rewrites \(\langle x,t\rangle-\langle y,t\rangle=\langle x-y,t\rangle\) with \texttt{inner\_\allowbreak{}sub\_\allowbreak{}left}, applies Cauchy–Schwarz \texttt{abs\_\allowbreak{}real\_\allowbreak{}inner\_\allowbreak{}le\_\allowbreak{}norm}, and commutes the product.

\end{enumerate}
\end{leanbox}

\section{\texorpdfstring{Characteristic differences controlled by the coupling distance}{Characteristic differences controlled by the coupling distance}}


\textbf{Claim.} Let \(H\) be a real inner product space with its Borel \(\sigma\)-algebra and a second countable topology, let \(\rho\) be a probability measure on \(H\times H\) with marginals \(\rho_1=(\mathrm{pr}_1)_\#\rho\) and \(\rho_2=(\mathrm{pr}_2)_\#\rho\), and assume \(\int\|x-y\|\,\rho(dx,dy)<\infty\). Then for every \(t\in H\),


\[\bigl|\varphi_{\rho_1}(t)-\varphi_{\rho_2}(t)\bigr|\le\|t\|\int\|x-y\|\,\rho(dx,dy).\]


\textbf{Proof.} By the change of variables formula for image measures,


\[\varphi_{\rho_1}(t)=\int e^{i\langle x,t\rangle}\,\rho(dx,dy),\qquad\varphi_{\rho_2}(t)=\int e^{i\langle y,t\rangle}\,\rho(dx,dy).\]


Both integrands are continuous, hence measurable, and bounded in modulus by \(1\); since \(\rho\) is a probability measure they are integrable. By linearity of the integral, the inequality \(|\int f\,d\rho|\le\int|f|\,d\rho\), the previous step and monotonicity of the integral,


\[\bigl|\varphi_{\rho_1}(t)-\varphi_{\rho_2}(t)\bigr|=\Bigl|\int\bigl(e^{i\langle x,t\rangle}-e^{i\langle y,t\rangle}\bigr)\,d\rho\Bigr|\le\int\bigl|e^{i\langle x,t\rangle}-e^{i\langle y,t\rangle}\bigr|\,d\rho\le\int\|t\|\,\|x-y\|\,d\rho .\]


The last integral is \(\|t\|\int\|x-y\|\,d\rho\), which is finite by the integrability hypothesis. \(\square\)

\begin{leanbox}{Lean: \texttt{charFun\_\allowbreak{}marginals\_\allowbreak{}norm\_\allowbreak{}sub\_\allowbreak{}le}}
\leansource{\texttt{charFun\_\allowbreak{}marginals\_\allowbreak{}norm\_\allowbreak{}sub\_\allowbreak{}le} \quad UnadjustedLangevin/\allowbreak{}CouplingLimit.lean, lines 42–64}
\begin{leancode}
\leanline{0}{theorem\ charFun\_marginals\_norm\_sub\_le}
\leanline{4}{(ρ\ :\ Measure\ (H\ ×\ H))\ [IsProbabilityMeasure\ ρ]}
\leanline{4}{(hint\ :\ Integrable\ (fun\ p\ :\ H\ ×\ H\ ↦\ ‖p.1\ -\ p.2‖)\ ρ)\ (t\ :\ H)\ :}
\leanline{4}{‖charFun\ (ρ.map\ Prod.fst)\ t\ -\ charFun\ (ρ.map\ Prod.snd)\ t‖\ ≤}
\leanline{6}{‖t‖\ *\ ∫\ p\ :\ H\ ×\ H,\ ‖p.1\ -\ p.2‖\ ∂ρ}
\end{leancode}
\leanlabel{How the Lean proof follows the argument}
In \texttt{char\allowbreak{}Fun\_\allowbreak{}marginals\_\allowbreak{}norm\_\allowbreak{}sub\_\allowbreak{}le}, \texttt{char\allowbreak{}Fun \ensuremath{\nu} t} is Mathlib's characteristic function \(\varphi_\nu(t)\), \texttt{\ensuremath{\rho}.\allowbreak{}map Prod.\allowbreak{}fst} and \texttt{\ensuremath{\rho}.\allowbreak{}map Prod.\allowbreak{}snd} are the marginals, and \texttt{hint} is the integrability of \(\|x-y\|\).



\begin{enumerate}


\item \texttt{hf} and \texttt{hg} are the integrability of the two characteristic tests under \(\rho\), proved by \texttt{Integrable.\allowbreak{}of\_\allowbreak{}bound} with the bound \(1\) from \texttt{Complex.\allowbreak{}norm\_\allowbreak{}exp\_\allowbreak{}of\allowbreak{}Real\_\allowbreak{}mul\_\allowbreak{}I} (\(|e^{i\theta}|=1\)).

\item The rewrite with \texttt{char\allowbreak{}Fun\_\allowbreak{}apply} and \texttt{integral\_\allowbreak{}map} (twice) expresses both characteristic functions as integrals over \(\rho\); \texttt{← integral\_\allowbreak{}sub hf hg} combines them into one integral.

\item The \texttt{calc} block is the displayed chain: \texttt{norm\_\allowbreak{}integral\_\allowbreak{}le\_\allowbreak{}integral\_\allowbreak{}norm}, then \texttt{integral\_\allowbreak{}mono} with the pointwise bound \texttt{characteristic\allowbreak{}Test\_\allowbreak{}norm\_\allowbreak{}sub\_\allowbreak{}le t p.\allowbreak{}1 p.\allowbreak{}2} and the integrability \texttt{hint.\allowbreak{}const\_\allowbreak{}mul \ensuremath{\Vert}t\ensuremath{\Vert}}, and finally \texttt{integral\_\allowbreak{}const\_\allowbreak{}mul}.

\end{enumerate}
\end{leanbox}

\section{\texorpdfstring{Continuity of the characteristic function in square mean}{Continuity of the characteristic function in square mean}}


\textbf{Claim.} Let \(P\) be a probability measure on \(\Omega\), \(H\) a real inner product space with its Borel \(\sigma\)-algebra, and \(t\in H\). Then the map \(X\mapsto\varphi_{P_X}(t)=\mathbb E\,e^{i\langle X,t\rangle}\) is continuous on \(L^2(P;H)\).




\textbf{Proof.} Let \(\ell(x)=\langle x,t\rangle\), a continuous linear functional \(H\to\mathbb R\) of norm at most \(\|t\|\). For \(X\in L^2(P;H)\), the real random variable \(\ell\circ X\) satisfies


\[\varphi_{P_X}(t)=\mathbb E\,e^{i\langle X,t\rangle}=\mathbb E\,e^{i\cdot1\cdot\ell(X)}=\varphi_{P_{\ell(X)}}(1),\qquad(2)\]


the characteristic function of the law of \(\ell(X)\) at frequency \(1\); this is the change of variables formula for the two image measures. The map \(X\mapsto\ell\circ X\) is a continuous linear map from \(L^2(P;H)\) to \(L^2(P;\mathbb R)\), since \(\|\ell\circ X\|_2\le\|t\|\,\|X\|_2\) and it does not depend on the representative of \(X\). The map \(Z\mapsto\varphi_{P_Z}(1)\) is continuous on \(L^2(P;\mathbb R)\) \cite[Section~3.5.7]{LangevinLean2026ULASC}; indeed, by (1), the Cauchy–Schwarz inequality and \(P(\Omega)=1\),


\[\bigl|\mathbb E\,e^{iZ}-\mathbb E\,e^{iZ'}\bigr|\le\mathbb E|Z-Z'|\le\|Z-Z'\|_2 .\]


By (2), \(X\mapsto\varphi_{P_X}(t)\) is the composition of these two continuous maps. \(\square\)

\begin{leanbox}{Lean: \texttt{continuous\_\allowbreak{}charFun\_\allowbreak{}map\_\allowbreak{}Lp\_\allowbreak{}hilbert}}
\leansource{\texttt{continuous\_\allowbreak{}charFun\_\allowbreak{}map\_\allowbreak{}Lp\_\allowbreak{}hilbert} \quad UnadjustedLangevin/\allowbreak{}CouplingLimit.lean, lines 71–84}
\begin{leancode}
\leanline{0}{theorem\ continuous\_charFun\_map\_Lp\_hilbert\ (t\ :\ H)\ :}
\leanline{4}{Continuous\ (fun\ X\ :\ Lp\ H\ 2\ P\ ↦\ charFun\ (P.map\ X)\ t)}
\end{leancode}
\leanlabel{How the Lean proof follows the argument}
\texttt{continuous\_\allowbreak{}char\allowbreak{}Fun\_\allowbreak{}map\_\allowbreak{}Lp\_\allowbreak{}hilbert} is stated without second countability (\texttt{omit [Second\allowbreak{}Countable\allowbreak{}Topology H] in}). In the proof, \texttt{L : H \ensuremath{\to}L[\ensuremath{\mathbb R}] \ensuremath{\mathbb R} := inner\allowbreak{}SL \ensuremath{\mathbb R} t} is the functional \(\ell\) as a continuous linear map.



\begin{enumerate}


\item \texttt{heq} is (2): Mathlib's \texttt{char\allowbreak{}Fun\_\allowbreak{}map\_\allowbreak{}eq\_\allowbreak{}char\allowbreak{}Fun\_\allowbreak{}map\_\allowbreak{}inner\_\allowbreak{}one} rewrites \(\varphi_{P_X}(t)\) as the characteristic function of \(\langle X,t\rangle\) at \(1\), and \texttt{Measure.\allowbreak{}map\_\allowbreak{}congr} with \texttt{L.\allowbreak{}coe\allowbreak{}Fn\_\allowbreak{}comp\allowbreak{}Lp X} identifies \(\langle X,t\rangle\) almost everywhere with the representative of the class \texttt{L.\allowbreak{}comp\allowbreak{}Lp X} \(=\ell\circ X\in L^2(P;\mathbb R)\) (the symmetry of the real inner product is \texttt{real\_\allowbreak{}inner\_\allowbreak{}comm}).

\item \texttt{hc} composes the scalar continuity theorem \texttt{continuous\_\allowbreak{}char\allowbreak{}Fun\_\allowbreak{}map\_\allowbreak{}Lp P (1 : \ensuremath{\mathbb R})} of the section on characteristic functions in \(L^2\) with the continuity of the bounded linear map \texttt{L.\allowbreak{}comp\allowbreak{}Lp\allowbreak{}L 2 P}; \texttt{convert hc using 1} and \texttt{funext X} reduce the goal to \texttt{heq}.

\end{enumerate}
\end{leanbox}

\section{\texorpdfstring{Equality of limiting laws from vanishing coupling distance}{Equality of limiting laws from vanishing coupling distance}}


Limits below are taken along a fixed nontrivial filter on an index set \(I\); the reader may take \(I=\mathbb N\) and limits as \(n\to\infty\), which is the case used later.




\textbf{Claim.} Let \(P\) be a probability measure on \(\Omega\) and \(H\) a real Hilbert space that is separable, with its Borel \(\sigma\)-algebra. Let \(X_i,X\in L^2(P;H)\) (\(i\in I\)), let \(\pi\) be a probability measure on \(H\), and let \(\rho_i\) be probability measures on \(H\times H\) such that for every \(i\)


\[(\mathrm{pr}_1)_\#\rho_i=P_{X_i},\qquad(\mathrm{pr}_2)_\#\rho_i=\pi,\qquad\int\|x-y\|\,\rho_i(dx,dy)<\infty,\]


and assume \(X_i\to X\) in \(L^2(P;H)\) and \(\int\|x-y\|\,\rho_i(dx,dy)\to0\). Then \(P_X=\pi\).




\textbf{Proof.} Fix \(t\in H\).




\emph{First limit.} By the previous step, \(X\mapsto\varphi_{P_X}(t)\) is continuous on \(L^2(P;H)\), so \(X_i\to X\) gives \(\varphi_{P_{X_i}}(t)\to\varphi_{P_X}(t)\).




\emph{Second limit.} By the second step applied to \(\rho_i\), whose marginals are \(P_{X_i}\) and \(\pi\),


\[0\le\bigl|\varphi_{P_{X_i}}(t)-\varphi_\pi(t)\bigr|\le\|t\|\int\|x-y\|\,\rho_i(dx,dy)\longrightarrow\|t\|\cdot0=0,\]


so, by the squeeze principle, \(\varphi_{P_{X_i}}(t)\to\varphi_\pi(t)\).




\emph{Conclusion.} Limits in \(\mathbb C\) along a nontrivial filter are unique, so \(\varphi_{P_X}(t)=\varphi_\pi(t)\). As \(t\) was arbitrary, \(P_X\) and \(\pi\) have the same characteristic function. Both are finite measures (\(P_X(H)=P(\Omega)=1\)) on a separable Hilbert space, and such measures are determined by their characteristic functions. Hence \(P_X=\pi\). \(\square\)




The couplings \(\rho_i\) need not be optimal, and only their first moments of \(\|x-y\|\) enter. Nontriviality of the filter excludes the degenerate situation in which every statement holds “eventually” and limits are not unique.

\begin{leanbox}{Lean: \texttt{measure\_\allowbreak{}map\_\allowbreak{}Lp\_\allowbreak{}eq\_\allowbreak{}of\_\allowbreak{}coupling\_\allowbreak{}cost\_\allowbreak{}tendsto}}
\leansource{\texttt{measure\_\allowbreak{}map\_\allowbreak{}Lp\_\allowbreak{}eq\_\allowbreak{}of\_\allowbreak{}coupling\_\allowbreak{}cost\_\allowbreak{}tendsto} \quad UnadjustedLangevin/\allowbreak{}CouplingLimit.lean, lines 92–114}
\begin{leancode}
\leanline{0}{theorem\ measure\_map\_Lp\_eq\_of\_coupling\_cost\_tendsto}
\leanline{4}{\{ι\ :\ Type*\}\ \{l\ :\ Filter\ ι\}\ [NeBot\ l]}
\leanline{4}{(Xn\ :\ ι\ →\ Lp\ H\ 2\ P)\ (X\ :\ Lp\ H\ 2\ P)}
\leanline{4}{(π\ :\ Measure\ H)\ [IsProbabilityMeasure\ π]}
\leanline{4}{(ρ\ :\ ι\ →\ Measure\ (H\ ×\ H))\ [∀\ i,\ IsProbabilityMeasure\ (ρ\ i)]}
\leanline{4}{(hX\ :\ Tendsto\ Xn\ l\ ({\leanfallback 𝓝}\ X))}
\leanline{4}{(hfst\ :\ ∀\ i,\ (ρ\ i).map\ Prod.fst\ =\ P.map\ (Xn\ i))}
\leanline{4}{(hsnd\ :\ ∀\ i,\ (ρ\ i).map\ Prod.snd\ =\ π)}
\leanline{4}{(hint\ :\ ∀\ i,\ Integrable\ (fun\ p\ :\ H\ ×\ H\ ↦\ ‖p.1\ -\ p.2‖)\ (ρ\ i))}
\leanline{4}{(hcost\ :\ Tendsto\ (fun\ i\ ↦\ ∫\ p\ :\ H\ ×\ H,\ ‖p.1\ -\ p.2‖\ ∂ρ\ i)\ l\ ({\leanfallback 𝓝}\ 0))\ :}
\leanline{4}{P.map\ X\ =\ π}
\end{leancode}
\leanlabel{How the Lean proof follows the argument}
\texttt{measure\_\allowbreak{}map\_\allowbreak{}Lp\_\allowbreak{}eq\_\allowbreak{}of\_\allowbreak{}coupling\_\allowbreak{}cost\_\allowbreak{}tendsto} works with a filter \texttt{l} on an index type \texttt{\ensuremath{\iota}} with \texttt{[Ne\allowbreak{}Bot l]} (nontriviality); \texttt{Tendsto f l ({\leanfallback 𝓝} a)} means that \(f(i)\to a\) along \texttt{l}. The hypotheses \texttt{h\allowbreak{}X}, \texttt{hfst}, \texttt{hsnd}, \texttt{hint}, \texttt{hcost} are those of the claim, and \texttt{[Complete\allowbreak{}Space H]} with the section variables give the Hilbert space setting.



\begin{enumerate}


\item \texttt{apply Measure.\allowbreak{}ext\_\allowbreak{}of\_\allowbreak{}char\allowbreak{}Fun} reduces the goal to equality of characteristic functions (Mathlib's uniqueness theorem for finite measures on a complete second-countable inner product space); \texttt{funext t} fixes the frequency.

\item \texttt{hactual} is the first limit, obtained by composing the continuity of the previous step at \(X\) with \texttt{h\allowbreak{}X}.

\item \texttt{hbound} is the coupling inequality \texttt{char\allowbreak{}Fun\_\allowbreak{}marginals\_\allowbreak{}norm\_\allowbreak{}sub\_\allowbreak{}le} with the marginals replaced via \texttt{hfst i} and \texttt{hsnd i}; \texttt{hnorm} applies \texttt{squeeze\_\allowbreak{}zero} with \texttt{hcost.\allowbreak{}const\_\allowbreak{}mul \ensuremath{\Vert}t\ensuremath{\Vert}}; \texttt{htarget} converts the vanishing norm of the difference into convergence (\texttt{tendsto\_\allowbreak{}iff\_\allowbreak{}norm\_\allowbreak{}sub\_\allowbreak{}tendsto\_\allowbreak{}zero}).

\item \texttt{tendsto\_\allowbreak{}nhds\_\allowbreak{}unique hactual htarget} is the uniqueness of limits.

\end{enumerate}
\end{leanbox}

\begin{thebibliography}{ULA-LSI}
\bibitem[ULA-SC]{LangevinLean2026ULASC}
LangevinLean Project. \emph{Unadjusted Langevin under Strong Convexity}. LL-01, 2026. Draft report, content revision d34ead06987c. \url{https://odakyildiz.github.io/LangevinLean/docs/pdf/UnadjustedLangevin/FullProof.pdf}.
\end{thebibliography}

\end{document}
